\subsection{COMPLEMENTS ON PLANE ROTATIONS}

In this no., V denotes a real vector space of dimension 2, provided with a scalar product (that is, a non-degenerate, positive, symmetric bilinear form) and an orientation. The measures of angles will be taken with respect to the base \(2\pi\) ; the principal measure of an angle between half-lines (resp. lines)

is thus a real number \(\theta\) such that \(0 \leqslant \theta < 2\pi\) (resp. \(0 \leqslant \theta < \pi\) ) (General Topology, Chap. VIII, §2, no. 3 and no. 6). By abuse of language, for any real number \(\theta\) , we shall use \(\theta\) to denote an angle whose measure is \(\theta\) and denote by \(\rho_{\theta}\) the rotation with angle \(\theta\) (Algebra, Chap. IX, §10, no. 3).

PROPOSITION 6. Let s be the orthogonal reflection with respect to a line D of V. If \(\Delta\) and \(\Delta'\) are two half-lines starting at the origin (resp. two lines passing through the origin) of V, we have

\[
(s (\widehat {\Delta}), s (\Delta^ {\prime})) \equiv - (\widehat {\Delta}, \Delta^ {\prime}) \pmod {2 \pi} \text {(resp. (mod. } \pi \text {))}.
\]

Let \(u\) be a rotation transforming \(\Delta\) into \(\Delta'\) . Since \(su\) is an orthogonal transformation of \(V\) of determinant \(-1\) , it is a reflection and thus \((su)^2 = 1\) . Consequently, \(u^{-1} = sus^{-1}\) transforms \(s(\Delta)\) into \(s(\Delta')\) , hence the proposition.

COROLLARY. Let D and \(D'\) be two lines of V and let \(\theta\) be a measure of the angle \((\widehat{\mathrm{D},\mathrm{D}'}\) ). Then \(s_{D'}s_{D} = \rho_{2\theta}\) .

We know that \(s_{D'} s_{D}\) is a rotation since it is of determinant 1. Let \(\Delta\) and \(\Delta'\) be two half-lines starting at the origin carried by D and \(D'\) . We have

\[
\begin{array}{r l} (\widehat {\Delta , s _ {\mathrm {D^ {\prime}}} s _ {\mathrm{D}}} (\Delta)) & \equiv (\widehat {\Delta , s _ {\mathrm {D^ {\prime}}}} (\Delta)) \equiv (\widehat {\Delta , \Delta^ {\prime}}) + (\Delta^ {\prime}, \widehat {s _ {\mathrm {D^ {\prime}}}} (\Delta)) \\ & \equiv (\widehat {\Delta , \Delta^ {\prime}}) + (s _ {\mathrm {D^ {\prime}}} (\widehat {\Delta^ {\prime}}), s _ {\mathrm {D^ {\prime}}} (\Delta)) \\ & \equiv (\widehat {\Delta , \Delta^ {\prime}}) - (\widehat {\Delta^ {\prime} , \Delta}) \equiv 2 (\widehat {\Delta , \Delta^ {\prime}}) \pmod {2 \pi}, \end{array}
\]

hence the corollary.

Now let \(\Delta\) and \(\Delta'\) be two half-lines of \(V\) such that

\[
\Delta \neq \Delta^ {\prime} \text { and } \Delta \neq - \Delta^ {\prime};
\]

and let \(s\) and \(s'\) be the orthogonal reflections with respect to the lines D and D' containing \(\Delta\) and \(\Delta'\) . Let \(\theta\) be the principal measure of the angle \((\widehat{\mathrm{D},\mathrm{D}'})\) . If \(\theta \in \pi\mathbf{Q}\) , denote by \(m\) the smallest integer \(\geqslant 1\) such that \(m\theta \in \pi\mathbf{Z}\) . If \(\theta \notin \pi\mathbf{Q}\) , put \(m = \infty\) . Let W be the group generated by \(s\) and \(s'\) .

PROPOSITION 7. The group W is dihedral (Chap. IV, §1, no. 2) of order \(2m\) . It consists of the rotations \(\rho_{2n\theta}\) and the products \(\rho_{2n\theta}s\) for \(n \in \mathbf{Z}\) . The transforms of D and D' by the elements of W are the transforms of D by the rotations \(\rho_{n\theta}\) for \(n \in \mathbf{Z}\) .

The Corollary to Prop. 6 shows that \(s's\) is of order \(m\) , which gives the first assertion. The elements of \(W\) are thus of the form \((s's)^n = \rho_{2n\theta}\) or \((s's)^n s = \rho_{2n\theta}s\) ; the last assertion follows from this, since \(D' = \rho_\theta(D)\) .

COROLLARY. Let C be the open angular sector formed by the union of the open half-lines \(\Delta_{1}\) starting at the origin and such that \(0 < (\widehat{\Delta, \Delta_{1}}) < \theta\) . Then no transform of D or \(\mathrm{D}'\) by an element of W meets C if and only if \(m\) is finite and

\[
\theta = \pi / m.
\]

If \(m = \infty\) , the image of the set \(n\theta\) ( \(n \in Z\) ) is dense in R/2πZ (General Topology, Chap. VII, §1, Prop. 11); the union of the transforms of D by the elements of W is thus dense in V and meets C. If m is finite and if \(\theta = k\pi/m\) with 1 \textless{} k \textless{} m, the integers k and m being relatively prime, there exists an integer h such that \(hk \equiv 1 \mod m\) ; then \((\widehat{\mathrm{D}, \rho_{h\theta}(\mathrm{D}))} \equiv \pi/m \pmod{\pi}\) , and \(\rho_{h\theta}(\mathrm{D})\) meets C. This shows that the condition is necessary. The converse is immediate.

Remark. Assume that \(m\) is finite and that \(\theta = \pi / m\) . If \(n \in \mathbf{Z}\) , let \(C_n\) be the union of the open half-lines \(\Delta_1\) starting at 0 such that

\[
n \theta <   (\widehat {\Delta , \Delta_ {1}}) <   (n + 1) \theta .
\]

The \(C_{n}\) for \(-m \leqslant n < m\) are connected open subsets forming a partition of \(E - \bigcup_{n} D_{n}\) (where \(D_{n} = \rho_{n\theta}(D)\) ). These are therefore the chambers determined in E by the system of m lines \(D_{n}\) ( \(1 \leqslant n \leqslant m\) ). We have \(C_{2k} = \rho_{2k\theta}(C)\) and \(C_{2k-1} = \rho_{2k\theta}s(C)\) . Moreover, \(C_{n} = C\) if and only if \(n \in 2mZ\) . Consequently, the group W permutes the chambers \(C_{n}\) simply-transitively.

We show finally that, if \(w \in W\) is such that the chambers C and \(w(\mathbf{C})\) are on opposite sides of the line D, then \(l(sw) = l(w) - 1\) (the lengths being taken with respect to \(S = \{s, s'\}\) ). Indeed, the assumption implies that \(w(\mathbf{C}) = \mathbf{C}_n\) with \(-m \leqslant n < 0\) . If n = -2k, we have \(w = (ss')^k\) and \(sw = s'(ss')^{k-1}\) , so \(l(w) = 2k\) and \(l(sw) = 2k - 1\) (Chap. IV, §1, no. 2, Remark). If \(n = -2k + 1\) , we have

\[
w = \left(s s ^ {\prime}\right) ^ {k - 1} s \quad \text { and } \quad s w = \left(s ^ {\prime} s\right) ^ {k - 1},
\]

hence \(l(w) = 2k - 1\) and \(l(sw) = 2k - 2\) . Q.E.D.

